\chapter{Complexity — Counter Prog Sim Run}
%%% generated by the blueprint pipeline from TCSlib.Complexity.SpaceComplexity.CounterProgSimRun
\section{Overview}

Declarations of the Lean module \texttt{TCSlib.Complexity.SpaceComplexity.CounterProgSimRun}.

\section{Declarations}

\begin{lemma}[Answer stability under appending bits]
\lean{Complexity.CPSim.ans_append}
\label{Complexity.CPSim.ans_append}
\leanok
\uses{Complexity.CPSim.ans}
\difficulty{4}
Let $O$ and $e$ be finite bit strings, and let $p<|O|$. The answer function at position
$p$ has the same value on $O$ and on the concatenation of $O$ with $e$, for either value
of its Boolean mode parameter.
\end{lemma}

\begin{lemma}[Answer beyond the output is false]
\lean{Complexity.CPSim.ans_of_le}
\label{Complexity.CPSim.ans_of_le}
\leanok
\uses{Complexity.CPSim.ans}
\difficulty{4}
Let $O$ be a finite bit string and let $p$ be a natural number with $|O|\le p$. The
answer function at position $p$ is false, for either value of its Boolean mode
parameter.
\end{lemma}

\begin{definition}[Bounded counter-program states]
\lean{Complexity.CPSim.Bd}
\label{Complexity.CPSim.Bd}
\leanok
\uses{Complexity.CounterProg.St}
For a counter-program state with $R$ registers, a natural-number bound $B$ holds when
every register value is at most $B$, the number of input symbols read is at most $B$,
and the length of the output string is at most $B$.
\end{definition}

\begin{lemma}[Counter Program Simulation Run]
\lean{Complexity.CPSim.sim_run}
\label{Complexity.CPSim.sim_run}
\leanok
\uses{Complexity.CPSim.A, Complexity.CPSim.Bd, Complexity.CPSim.G, Complexity.CPSim.ans, Complexity.CPSim.ans_append, Complexity.CPSim.ans_of_le, Complexity.CPSim.cf, Complexity.CPSim.ev, Complexity.CPSim.sim_step, Complexity.CounterProg.St, Complexity.CounterProg.run, Complexity.CounterProg.run_of_halted, Complexity.CounterProg.run_out, Complexity.CounterProg.run_succ, Complexity.CounterProg.run_zero, Complexity.CounterProg.step, Complexity.LogProg.AHalt, Complexity.LogProg.AReach.halt, Turing.pairEncode}
\difficulty{6}
Let $R$ be a number of registers, let $\Lambda$ be a set of labels, and let $P$ be a
counter program on these registers and labels. Fix a starting label $l_0$, a Boolean
mode $lm$, natural numbers $n,p,B$, an input word $u$, and two Boolean oracles $o_0,o_1$
on bit strings. Suppose that for every $h\in\mathbb N$, the oracle answers on the
encoding of $(1^n,\mathrm{bits}(h))$ are $o_0=1$ exactly when $h<|u|$ and $o_1=u_h$,
with out-of-range bits interpreted as $0$. If a counter-program state $s$ is at label
$l$, its output has length at most $p$, its time register $tm$ is at most $B$, and after
$k$ steps on input $u$ it is halted, assume also that at every time $j\le k$ all its
registers, its input position, and its output length are at most $B$. Then the auxiliary
program started in its simulation state for $l$, with register values encoding those of
$s$, its input position, output length, and $tm$, halts while satisfying its prescribed
invariant at every earlier step, and returns the bit determined from the final output by
the rule that when $lm$ is true it tests whether position $p$ is within the output, and
otherwise reads output position $p$ with default $0$.
\end{lemma}

\begin{theorem}[Bound on output length]
\lean{Complexity.CounterProg.length_out_le}
\label{Complexity.CounterProg.length_out_le}
\leanok
\uses{Complexity.CounterProg.Instr, Complexity.CounterProg.init, Complexity.CounterProg.run, Complexity.CounterProg.run_init_out_le}
\difficulty{3}
Let $P$ be a counter program with $R$ registers and labels $Lambda$, let
$l_0\inLambda$, and let $x$ be a binary input. For natural numbers $C,c,n,t$, if $t\le
C(n+1)^c$, then the output after running $P$ for $t$ steps from its initial state at
$l_0$ has length at most $(C+1)^2(n+1)^{2c}$.
\end{theorem}

\begin{theorem}[Counter programs preserve unary logspace]
\lean{Complexity.UnaryLogspace.counterProg}
\label{Complexity.UnaryLogspace.counterProg}
\leanok
\uses{Complexity.CPSim.A, Complexity.CPSim.Bd, Complexity.CPSim.ans, Complexity.CPSim.cf, Complexity.CPSim.ev, Complexity.CPSim.preS_A, Complexity.CPSim.sim_run, Complexity.CounterProg.Instr, Complexity.CounterProg.St, Complexity.CounterProg.init, Complexity.CounterProg.run, Complexity.CounterProg.run_init_out_le, Complexity.CounterProg.run_pos_le, Complexity.CounterProg.run_regs_le, Complexity.LOGSPACE, Complexity.LogProg.AHalt.mono, Complexity.LogProg.AHalt.step, Complexity.LogProg.PreS, Complexity.LogProg.ValidPlain, Complexity.LogProg.arm_decides_poly, Complexity.LogProg.arun, Complexity.LogProg.astep, Complexity.LogProg.bitsVal, Complexity.LogProg.canon_bits, Complexity.LogProg.canon_eq_bits, Complexity.LogProg.dbl_replicate, Complexity.LogProg.indicator_eq_decide, Complexity.UnaryLogspace, Complexity.mem_uBit, Complexity.mem_uLen, Complexity.uBit, Complexity.uLen, Complexity.validPlain_of_mem_unaryIdx, Turing.MultiTapeTM.indicator, Turing.pairEncode, Turing.pairEncode_eq_dbl}
\difficulty{7}
Let $R\in\mathbb{N}$, let $\Lambda$ be a finite set with decidable equality, and let $P$
be a counter program with $R$ registers and labels in $\Lambda$, with initial label
$\ell_0$. Suppose $u,g:\mathbb{N}\to\{0,1\}^*$, and both the bit-membership and
output-length languages associated to $u$ are decidable in logarithmic space. If there
are natural numbers $C,c$ such that, for every $n$, the program on input $u(n)$ halts
within $C(n+1)^c$ steps and has output $g(n)$, then both corresponding languages for $g$
are decidable in logarithmic space.
\end{theorem}
